\documentclass[10pt,a4paper]{amsart}
\usepackage{amsmath,amssymb,amsthm,graphicx}
\usepackage{verbatim}
\newenvironment{reference}{\comment}{\endcomment}
%\newenvironment{reference}{}{}
\newenvironment{slogan}{\comment}{\endcomment}
\newenvironment{history}{\comment}{\endcomment}
% For commutative diagrams we use Xy-pic
\usepackage[all]{xy}

% We use 2cell for 2-commutative diagrams.
\xyoption{2cell}
% Theorem environments.
%
\theoremstyle{plain}
\newtheorem{theorem}[subsection]{Theorem}
\newtheorem{proposition}[subsection]{Proposition}
\newtheorem{lemma}[subsection]{Lemma}

\theoremstyle{definition}
\newtheorem{definition}[subsection]{Definition}
\newtheorem{example}[subsection]{Example}
\newtheorem{exercise}[subsection]{Exercise}
\newtheorem{situation}[subsection]{Situation}

\theoremstyle{remark}
\newtheorem{remark}[subsection]{Remark}
\newtheorem{remarks}[subsection]{Remarks}

\numberwithin{equation}{subsection}

% Macros
%
\def\lim{\mathop{\mathrm{lim}}\nolimits}
\def\colim{\mathop{\mathrm{colim}}\nolimits}
\def\Spec{\mathop{\mathrm{Spec}}}
\def\Hom{\mathop{\mathrm{Hom}}\nolimits}
\def\Ext{\mathop{\mathrm{Ext}}\nolimits}
\def\SheafHom{\mathop{\mathcal{H}\!\mathit{om}}\nolimits}
\def\SheafExt{\mathop{\mathcal{E}\!\mathit{xt}}\nolimits}
\def\Sch{\mathit{Sch}}
\def\Mor{\mathop{\mathrm{Mor}}\nolimits}
\def\Ob{\mathop{\mathrm{Ob}}\nolimits}
\def\Sh{\mathop{\mathit{Sh}}\nolimits}
\def\NL{\mathop{N\!L}\nolimits}
\def\CH{\mathop{\mathrm{CH}}\nolimits}
\def\proetale{{pro\text{-}\acute{e}tale}}
\def\etale{{\acute{e}tale}}
\def\QCoh{\mathit{QCoh}}
\def\Ker{\mathop{\mathrm{Ker}}}
\def\Im{\mathop{\mathrm{Im}}}
\def\Coker{\mathop{\mathrm{Coker}}}
\def\Coim{\mathop{\mathrm{Coim}}}

% Boxtimes
%
\DeclareMathSymbol{\boxtimes}{\mathbin}{AMSa}{"02}

%
% Macros for moduli stacks/spaces
%
\def\QCohstack{\mathcal{QC}\!\mathit{oh}}
\def\Cohstack{\mathcal{C}\!\mathit{oh}}
\def\Spacesstack{\mathcal{S}\!\mathit{paces}}
\def\Quotfunctor{\mathrm{Quot}}
\def\Hilbfunctor{\mathrm{Hilb}}
\def\Curvesstack{\mathcal{C}\!\mathit{urves}}
\def\Polarizedstack{\mathcal{P}\!\mathit{olarized}}
\def\Complexesstack{\mathcal{C}\!\mathit{omplexes}}
% \Pic is the operator that assigns to X its picard group, usage \Pic(X)
% \Picardstack_{X/B} denotes the Picard stack of X over B
% \Picardfunctor_{X/B} denotes the Picard functor of X over B
\def\Pic{\mathop{\mathrm{Pic}}\nolimits}
\def\Picardstack{\mathcal{P}\!\mathit{ic}}
\def\Picardfunctor{\mathrm{Pic}}
\def\Deformationcategory{\mathcal{D}\!\mathit{ef}}
\usepackage{cleveref}
\crefname{originallemma}{original lemma}{original lemmas}
\Crefname{originallemma}{Original lemma}{Original lemmas}
\makeatletter
\newlabel{lemma-compute-for-effective-Cartier-algebraic}{{lemma-compute-for-effective-Cartier-algebraic}{source}{}{sourcepointer-1}{}}
\newlabel{lemma-compute-for-effective-Cartier-algebraic@cref}{{[originallemma][0][]lemma-compute-for-effective-Cartier-algebraic}{[1][0][]source}{}{}{}}
\makeatother
\crefname{originalsection}{original section}{original sections}
\Crefname{originalsection}{Original section}{Original sections}
\makeatletter
\newlabel{section-trivial}{{section-trivial}{source}{}{sourcepointer-2}{}}
\newlabel{section-trivial@cref}{{[originalsection][0][]section-trivial}{[1][0][]source}{}{}{}}
\makeatother
\title{Independence of the presentation in the construction of the twisted inverse image functor}
\author{The Stacks Project Authors}
\date{stacks-project:dualizing.tex:local-snapshot-20261012; GNU FDL 1.2 or later}
\begin{document}
\maketitle
\noindent\textit{Editorial note.} All delivered mathematical text is the author's own, copied from the pinned Stacks Project chapter source sources/stacks-dualizing/pinned-original-source.tex; no statement or proof step was rewritten, reordered or removed.\par\medskip
\noindent\textit{Editorial note.} Preamble provenance: the Stacks Project keeps one single shared preamble for every chapter; this card takes it directly from the staged sources/stacks-dualizing/preamble.tex as three verbatim spans (lines 8-12, 14-18 and 154-217), so no derived file was written and nothing was added or redefined. The wrapper class and the base packages are supplied by the assembler.\par\medskip
\noindent\textit{Editorial note.} Bibliography provenance: neither the counted statement nor its proof contains a citation (this lane rescans the whole span before writing the card), so the card carries no bibliography and no entry was reconstructed.\par\medskip
\noindent\textit{Editorial note.} Selection provenance: this card was opened under the file-based control contract (CONTROL.md), with the mechanical pre-card gate run on its statement and proof interval before the card was written.\par\medskip
\noindent\textit{Allowed standard prerequisites (editorial).}\begin{itemize}
\item The identification of the functor attached to a presentation as the right adjoint of the restriction along it. Invoked inside the counted proof (source lines 4612-4615). Makes the functor for the composed presentation the composite of the two functors.
\item The computation of the functor for the extension of a presentation by a single variable. Invoked at the end of the counted proof (source lines 4624-4625). Identifies the single step with the shift of the tensor product.
\end{itemize}\medskip
\ifdefined\hypertarget\hypertarget{sourcepointer-1}{}\fi
\noindent\textit{Source reference.} Original lemma lemma-compute-for-effective-Cartier-algebraic: sources/stacks-dualizing/pinned-original-source.tex, source lines 2621--2621. Navigation pointer only: the counted text refers to the Stacks item lemma-compute-for-effective-Cartier-algebraic, which is an existing labeled result of this same chapter and is not part of this card. The pointer carries the target's own label as its number and its definition line as its span; its content is not reproduced here.\par\medskip
\ifdefined\hypertarget\hypertarget{sourcepointer-2}{}\fi
\noindent\textit{Source reference.} Original section section-trivial: sources/stacks-dualizing/pinned-original-source.tex, source lines 2416--2416. Navigation pointer only: the counted text refers to the Stacks item section-trivial, which is an existing labeled result of this same chapter and is not part of this card. The pointer carries the target's own label as its number and its definition line as its span; its content is not reproduced here.\par\medskip
\begin{lemma}
\label{lemma-well-defined}
Let $\varphi : R \to A$ be a finite type homomorphism of
Noetherian rings. The functor $\varphi^!$ is well defined
up to isomorphism.
\end{lemma}

\begin{proof}
Suppose that $\psi_1 : P_1 = R[x_1, \ldots, x_n] \to A$ and
$\psi_2 : P_2 = R[y_1, \ldots, y_m] \to A$ are two
surjections from polynomial rings onto $A$. Then we get a
commutative diagram
$$
\xymatrix{
R[x_1, \ldots, x_n, y_1, \ldots, y_m]
\ar[d]^{x_i \mapsto g_i} \ar[rr]_-{y_j \mapsto f_j} & &
R[x_1, \ldots, x_n] \ar[d] \\
R[y_1, \ldots, y_m] \ar[rr] & & A
}
$$
where $f_j$ and $g_i$ are chosen such that $\psi_1(f_j) = \psi_2(y_j)$
and $\psi_2(g_i) = \psi_1(x_i)$. By symmetry it suffices to prove
the functors defined using $P \to A$ and $P[y_1, \ldots, y_m] \to A$
are isomorphic. By induction we may assume $m = 1$. This reduces
us to the case discussed in the next paragraph.

\medskip\noindent
Here $\psi : P \to A$ is given and $\chi : P[y] \to A$ induces
$\psi$ on $P$. Write $Q = P[y]$.
Choose $g \in P$ with $\psi(g) = \chi(y)$.
Denote $\pi : Q \to P$ the $P$-algebra map
with $\pi(y) = g$. Then $\chi = \psi \circ \pi$ and hence
$\chi^! = \psi^! \circ \pi^!$ as both are
adjoint to the restriction functor $D(A) \to D(Q)$ by the material
in Section \ref{section-trivial}. Thus
$$
\chi^!\left(K \otimes_R^\mathbf{L} Q\right)[n + 1] =
\psi^!\left(\pi^!\left(K \otimes_R^\mathbf{L} Q\right)[1]\right)[n]
$$
Hence it suffices to show that
$\pi^!(K \otimes_R^\mathbf{L} Q[1]) = K \otimes_R^\mathbf{L} P$
Thus it suffices to show that the functor
$\pi^!(-) : D(Q) \to D(P)$
is isomorphic to $K \mapsto K \otimes_Q^\mathbf{L} P[-1]$.
This follows from Lemma \ref{lemma-compute-for-effective-Cartier-algebraic}.
\end{proof}

\end{document}
